\section{Rastreabilidade e material suplementar}
\label{ap:fontes}
\begin{singlespace}

O projeto que acompanha este relatório contém os logs em \codigo{dados/logs/}, os CSV derivados em \codigo{dados/}, os gráficos vetoriais e rasterizados em \codigo{imagens/} e as fontes bibliográficas em \codigo{references.bib}. A Tabela~\ref{tab:fontes} permite localizar cada ensaio. O inventário armazena o SHA-256 integral de cada arquivo.

\begin{singlespace}
\begingroup
\footnotesize
\setlength{\tabcolsep}{3pt}
\renewcommand{\arraystretch}{1.18}
\begin{longtable}{@{}lP{14cm}@{}}
\caption{Correspondência entre IDs e arquivos de coleta.}\label{tab:fontes}\\
\toprule
ID & Arquivo em dados/logs/\\
\midrule
\endfirsthead
\multicolumn{2}{c}{\tablename\ \thetable\ -- continuação}\\
\toprule
ID & Arquivo em dados/logs/\\
\midrule
\endhead
\midrule\multicolumn{2}{r}{Continua na próxima página.}\\
\endfoot
\bottomrule
\endlastfoot
01 & \path{01_CUSTOM_PREEMPTIVO_80MHz_UNICORE.txt}\\
02 & \path{02_CUSTOM_PREEMPTIVO_240MHz_UNICORE.txt}\\
03 & \path{03_CUSTOM_COOPERATIVO_80MHz_UNICORE.txt}\\
04 & \path{04_CUSTOM_COOPERATIVO_240MHz_UNICORE.txt}\\
05 & \path{05_DM_PREEMPTIVO_80MHz_UNICORE.txt}\\
06 & \path{06_DM_PREEMPTIVO_240MHz_UNICORE.txt}\\
07 & \path{07_DM_COOPERATIVO_80MHz_UNICORE.txt}\\
08 & \path{08_DM_COOPERATIVO_240MHz_UNICORE.txt}\\
E01 & \path{EXTRA_01_CUSTOM_PREEMPTIVO_160MHz_UNICORE.txt}\\
E02 & \path{EXTRA_02_CUSTOM_COOPERATIVO_160MHz_UNICORE.txt}\\
E03 & \path{EXTRA_03_DM_PREEMPTIVO_160MHz_UNICORE.txt}\\
E04 & \path{EXTRA_04_DM_COOPERATIVO_160MHz_UNICORE.txt}\\
E05 & \path{EXTRA_05_DM_PREEMPTIVO_80MHz_UNICORE_SLICING_OFF.txt}\\
E06 & \path{EXTRA_06_DM_PREEMPTIVO_240MHz_UNICORE_SLICING_OFF.txt}\\
E07 & \path{EXTRA_07_RM_PREEMPTIVO_240MHz_UNICORE_SLICING_ON.txt}\\
E08 & \path{EXTRA_08_RM_PREEMPTIVO_80MHz_UNICORE_SLICING_ON.txt}\\
E09 & \path{EXTRA_09_RM_COOPERATIVO_80MHz_UNICORE.txt}\\
E10 & \path{EXTRA_10_RM_COOPERATIVO_240MHz_UNICORE.txt}\\
\end{longtable}
\endgroup
\noindent{\footnotesize Fonte: inventário das coletas; hashes integrais em dados/inventario.csv.}
\end{singlespace}


A cópia de \codigo{codigo/main.c} representa o código consultado durante a análise; \codigo{codigo/sdkconfig_consultado.txt} representa a configuração disponível nessa consulta. Ambos auxiliam a auditoria, mas não identificam retroativamente cada binário das coletas. O SDK local consultado corresponde à revisão \codigo{fff9895c82d744c7237be8847347bdd1b07c6643}. Os trechos do driver/macro usados para interpretar o yield na ISR foram preservados em \codigo{fontes/trechos_sdk_consultado.txt}.

O enunciado completo está em \codigo{fontes/enunciado.pdf}, e a cópia local do exemplo inicial do professor em \codigo{fontes/Drone_exemplo_oTask_referencia.txt}. O endereço do repositório original é citado nas referências~\cite{viel_codigo}. Os documentos online do ESP-IDF identificados como \textit{latest} podem evoluir; o exemplo touch foi referenciado na revisão do SDK consultado~\cite{idf_touch_example}.

As tabelas podem ser regeneradas a partir dos CSV incluídos executando \codigo{python ferramentas/gerar_tabelas.py}. Essa etapa é externa ao firmware e não interfere nas medições. O pacote contém somente as observações existentes e suas derivações; não inclui ensaios sintéticos para preencher combinações não executadas.
\end{singlespace}
